\chapter{วิธีดำเนินงานวิจัย}
การวิจัยเรื่อง ``การพัฒนาทักษะเกษตรกรอัจฉริยะและการจำแนกกลุ่มเป้าหมายด้วยเทคโนโลยีการเรียนรู้ของเครื่อง'' 
เป็นการวิจัยแบบผสมผสาน (Mixed Method Research) เพื่อศึกษาผลกระทบของโครงการและพัฒนาแบบจำลองจำแนกกลุ่มเกษตรกร 
ผู้วิจัยได้กำหนดขั้นตอนการดำเนินงานวิจัยตามลำดับ ดังนี้

\section{รูปแบบการวิจัย (Research Design)}
การวิจัยเรื่อง “การพัฒนาทักษะเกษตรกรอัจฉริยะและการจำแนกกลุ่มเป้าหมายด้วยเทคโนโลยีการเรียนรู้ของเครื่อง” 
เป็นการวิจัยแบบผสมผสาน (Mixed Method Research) โดยบูรณาการการวิเคราะห์เชิงเศรษฐมิติร่วมกับเทคนิคการเรียนรู้ของเครื่อง (Machine Learning) 
ในการดำเนินการวิจัย ผู้วิจัยได้แบ่งกระบวนการออกเป็น 2 ส่วนหลัก ได้แก่ การวิจัยเชิงปริมาณและการวิจัยเชิงคุณภาพ

\subsection{การวิจัยเชิงปริมาณ (Quantitative Research)}

ใช้ข้อมูลทุติยภูมิในการประเมินผลกระทบของโครงการด้วยวิธีทางเศรษฐมิติ ได้แก่ 
Propensity Score Matching (PSM) และ Difference-in-Differences (DID) 
รวมทั้งประยุกต์ใช้เทคนิคการเรียนรู้ของเครื่อง (Machine Learning) 
เพื่อวิเคราะห์ความแตกต่างของผลลัพธ์ในระดับรายบุคคล

\subsection{การวิจัยเชิงคุณภาพ (Qualitative Research)}

ใช้การสัมภาษณ์เชิงลึกเพื่อเสริมความเข้าใจเชิงบริบท 
และสนับสนุนการตีความผลการวิเคราะห์เชิงปริมาณให้มีความรอบด้าน
% -----------------------

\section{ประชากรและกลุ่มตัวอย่าง (Population and Sample)}

\subsection{ประชากร( Population)}  
ประชากรที่ใช้ในการวิจัย ได้แก่ เกษตรกรที่ปรากฏอยู่ในฐานข้อมูลโครงการ Smart Farmer ของกรมวิชาการเกษตร ในช่วงปี พ.ศ. 2565–2566

\subsection{กลุ่มตัวอย่าง (Sample)}  
กำหนดขนาดกลุ่มตัวอย่างประมาณ 350–400 ราย โดยจำแนกเป็น 2 กลุ่ม ได้แก่

\begin{itemize}
\item กลุ่มทดลอง (Treatment Group): เกษตรกรที่เข้าร่วมโครงการ Smart Farmer
\item กลุ่มควบคุม (Control Group): เกษตรกรที่ไม่เข้าร่วมโครงการ
\end{itemize}

\subsection{วิธีการคัดเลือกกลุมตัวอย่าง}  
การคัดเลือกกลุ่มตัวอย่างใช้วิธีการสุ่มแบบแบ่งชั้นภูมิ (Stratified Random Sampling) เพื่อให้กลุ่มตัวอย่างครอบคลุมเกษตรกรในหลากหลายพื้นที่โดยกำหนดชั้นภูมิตามพื้นที่/จังหวัดหรือเขตการผลิตตามข้อมูลที่มีอยู่ จากนั้นจึงสุ่มตัวอย่างภายในแต่ละชั้นภูมิ

% -----------------------

\section{เครื่องมือที่ใช้ในการวิจัย}

ผู้วิจัยเลือกใช้เครื่องมือในการเก็บรวบรวมและวิเคราะห์ข้อมูล ดังนี้

\begin{itemize}

\item \textbf{แบบสัมภาษณ์เชิงลึก (In-depth Interview Form)} 
สำหรับสัมภาษณ์เกษตรกรและผู้เชี่ยวชาญเพื่อเจาะลึกพฤติกรรมและปัญหา

\item \textbf{แบบประเมินทักษะเกษตรกร (Skill Assessment Form)} 
ใช้ประเมินระดับความรู้ความสามารถก่อนและหลังเข้าร่วมโครงการ

\item \textbf{แบบบันทึกข้อมูลการฝึกอบรม} 
รวบรวมประวัติหลักสูตร ระยะเวลา และรูปแบบกิจกรรมที่เกษตรกรเข้าร่วม

\item \textbf{โมเดลการเรียนรู้ของเครื่อง (Machine Learning Model)} 
อัลกอริทึมสำหรับวิเคราะห์และจำแนกกลุ่มเกษตรกร (Clustering \& Classification)

\end{itemize}

% -----------------------

\section{การเก็บรวบรวมข้อมูล (Data Collection)}

\subsection{ข้อมูลทุติยภูมิ (Secondary Data)}

รวบรวมจากฐานข้อมูลภาครัฐและฐานข้อมูลโครงการ Smart Farmer 
ประกอบด้วยข้อมูลพื้นฐานของเกษตรกร ประวัติการเข้าร่วมโครงการหรือการฝึกอบรม 
และคะแนนทักษะก่อนและหลังการอบรม (Pre-test/Post-test)

\subsection{ข้อมูลปฐมภูมิ (Primary Data)}

ดำเนินการเก็บข้อมูลเพิ่มเติมเพื่อเสริมความเข้าใจเชิงบริบท 
และใช้ประกอบการอภิปรายผล โดยใช้วิธีการสัมภาษณ์เชิงลึก (In-depth Interview) 
และการสังเกตการณ์ภาคสนาม (Field Observation)

% -----------------------

\section{ตัวแปรที่ใช้ในการศึกษา (Variables)}

การศึกษานี้กำหนดตัวแปรตามและตัวแปรอิสระเพื่อใช้ในการประเมินผลกระทบของโครงการ 
และการวิเคราะห์ด้วยเทคนิคการเรียนรู้ของเครื่อง โดยสรุปได้ดังนี้

\subsection{ตัวแปรตาม (Dependent Variable)}

\begin{itemize}
\item การเปลี่ยนแปลงของทักษะเกษตรกร (Skill Change) 
วัดจากผลต่างคะแนนทักษะหลังเข้าร่วมโครงการและก่อนเข้าร่วมโครงการ
\end{itemize}

\subsection{ตัวแปรอิสระ (Independent Variables)}

\begin{itemize}
\item \textbf{Human Capital}: อายุ ระดับการศึกษา และประสบการณ์ด้านการเกษตร

\item \textbf{Physical Capital}: ขนาดที่ดิน และการเข้าถึงระบบชลประทาน

\item \textbf{Financial and Risk Factors}: ภาระหนี้ รายได้อื่น และความเสี่ยง

\item \textbf{Social Capital}: การเป็นสมาชิกกลุ่ม และการเข้าถึงข้อมูลหรือช่องทางการเรียนรู้

\item \textbf{Baseline Skills}: ระดับทักษะก่อนเข้าร่วมโครงการ
\end{itemize}

% -----------------------

\section{การวิเคราะห์ผลกระทบของโครงการด้วย PSM-DID}

เพื่อประเมินผลกระทบของโครงการ Smart Farmer ต่อการเปลี่ยนแปลงทักษะของเกษตรกร 
ผู้วิจัยประยุกต์ใช้วิธี Propensity Score Matching (PSM) ร่วมกับ 
Difference-in-Differences (DID) เพื่อลดอคติจากการคัดเลือกเข้าร่วมโครงการ 
และเพื่อประมาณผลกระทบเชิงสาเหตุได้อย่างเหมาะสม

% ----------------------

\subsection{Propensity Score Matching (PSM)}

ใช้เพื่อจับคู่เกษตรกรในกลุ่มทดลองและกลุ่มควบคุมที่มีคุณลักษณะใกล้เคียงกัน 
บนพื้นฐานของตัวแปรก่อนเข้าร่วมโครงการ เพื่อเพิ่มความสามารถในการเปรียบเทียบ (comparability)

% ----------------------

\subsection{Difference-in-Differences (DID)}

วิธี Difference-in-Differences (DID) ใช้เพื่อประมาณผลกระทบสุทธิของโครงการ 
โดยเปรียบเทียบการเปลี่ยนแปลงของคะแนนทักษะระหว่างกลุ่มทดลองและกลุ่มควบคุม 
ในช่วงเวลาก่อนและหลังการดำเนินโครงการ โดยสมการสามารถเขียนได้ดังนี้

\section{การวิเคราะห์ผลกระทบของโครงการด้วย PSM-DID}

เพื่อประเมินผลกระทบของโครงการ Smart Farmer ต่อการเปลี่ยนแปลงทักษะของเกษตรกร 
ผู้วิจัยประยุกต์ใช้วิธี Propensity Score Matching (PSM) ร่วมกับ 
Difference-in-Differences (DID) เพื่อลดอคติจากการคัดเลือกเข้าร่วมโครงการ 
และเพื่อประมาณผลกระทบเชิงสาเหตุได้อย่างเหมาะสม

% ----------------------

\subsection{Propensity Score Matching (PSM)}

ใช้เพื่อจับคู่เกษตรกรในกลุ่มทดลองและกลุ่มควบคุมที่มีคุณลักษณะใกล้เคียงกัน 
บนพื้นฐานของตัวแปรก่อนเข้าร่วมโครงการ เพื่อเพิ่มความสามารถในการเปรียบเทียบ (comparability)

% ----------------------

\subsection{Difference-in-Differences (DID)}

วิธี Difference-in-Differences (DID) ใช้เพื่อประมาณผลกระทบสุทธิของโครงการ 
โดยเปรียบเทียบการเปลี่ยนแปลงของคะแนนทักษะระหว่างกลุ่มทดลองและกลุ่มควบคุม 
ในช่วงเวลาก่อนและหลังการดำเนินโครงการ โดยสมการสามารถเขียนได้ดังนี้

{\fontsize{12}{15}\selectfont

\begin{equation}
Skill_{it} = \beta_0 + \beta_1 Treatment_i + \beta_2 Time_t 
+ \beta_3 (Treatment_i \times Time_t) + \epsilon_{it}
\end{equation}

โดยที่

\begin{itemize}
\item $Treatment_i$ หมายถึง สถานะการเข้าร่วมโครงการ (1 = เข้าร่วม, 0 = ไม่เข้าร่วม)
\item $Time_t$ หมายถึง ช่วงเวลา (1 = หลังโครงการ, 0 = ก่อนโครงการ)
\item $\beta_3$ แสดงผลกระทบสุทธิของโครงการต่อการเปลี่ยนแปลงทักษะ
\end{itemize}

}
% ----------------------

\subsection{การวิเคราะห์เชิงพรรณนา (Descriptive Statistics)}

ใช้ค่าสถิติพื้นฐาน ได้แก่ ค่าร้อยละ ค่าเฉลี่ย และส่วนเบี่ยงเบนมาตรฐาน 
เพื่ออธิบายลักษณะทั่วไปของกลุ่มตัวอย่าง

% ----------------------

\subsection{การวิเคราะห์ด้วยแบบจำลองทางเศรษฐมิติ}

ใช้แบบจำลองการถดถอยข้อมูลแบบพาแนล (Panel Data Regression Model) 
เพื่อวิเคราะห์ความสัมพันธ์ระหว่างการเข้าร่วมโครงการและการเปลี่ยนแปลงทักษะ 
โดยอาศัยแนวคิดของ Wooldridge (2009)

% -----------------------

\section{การวิเคราะห์ด้วย Machine Learning}

เพื่อวิเคราะห์ความแตกต่างของผลลัพธ์ในระดับรายบุคคลและพัฒนาแบบจำลองสำหรับการทำนายและการจำแนกกลุ่มเกษตรกร ผู้วิจัยประยุกต์ใช้เทคนิคการเรียนรู้ของเครื่อง (Machine Learning) ตามขั้นตอนการวิเคราะห์ ดังนี้

\subsection{การเตรียมข้อมูล (Data Preprocessing)}

ก่อนนำข้อมูลเข้าสู่การสร้างแบบจำลอง ผู้วิจัยดำเนินการเตรียมข้อมูลให้เหมาะสมกับอัลกอริทึม ได้แก่ การทำความสะอาดข้อมูล การแปลงข้อมูล และการปรับสเกลข้อมูล

\begin{itemize}
\item \textbf{การทำความสะอาดข้อมูล (Data Cleaning)}: ตรวจสอบและจัดการข้อมูลสูญหาย (Missing Values) และค่าผิดปกติ (Outliers) ที่อาจส่งผลต่อความน่าเชื่อถือของการประมาณค่า

\item \textbf{การแปลงข้อมูล (Data Transformation)}: แปลงตัวแปรเชิงคุณภาพให้เป็นตัวแปรเชิงปริมาณตามความเหมาะสม เช่น การใช้ One-Hot Encoding หรือ Label Encoding เพื่อให้สามารถนำไปใช้ในการสร้างแบบจำลองได้

\item \textbf{การปรับสเกลข้อมูล (Feature Scaling)}: ปรับมาตรฐานของตัวแปรที่มีหน่วยวัดแตกต่างกันให้อยู่ในระดับที่เปรียบเทียบได้ เช่น Standardization หรือ Normalization เพื่อสนับสนุนประสิทธิภาพของอัลกอริทึม
\end{itemize}

\subsection{การคัดเลือกตัวแปร (Feature Selection)}

ผู้วิจัยใช้การวิเคราะห์สหสัมพันธ์ (Correlation Analysis) เพื่อตรวจสอบความสัมพันธ์ระหว่างตัวแปร และใช้แนวคิด Feature Importance เช่น จากแบบจำลองเชิงต้นไม้ เพื่อคัดเลือกตัวแปรสำคัญ ลดมิติข้อมูล และเพิ่มประสิทธิภาพของโมเดล

\subsection{การพัฒนาโมเดล (Model Development)}

ผู้วิจัยพัฒนาโมเดลด้วยการเรียนรู้แบบมีผู้สอน (Supervised Learning) สำหรับการทำนายและจำแนกผลลัพธ์ โดยเปรียบเทียบประสิทธิภาพของอัลกอริทึม ได้แก่ Decision Tree, Random Forest และ Support Vector Machine (SVM)

\begin{itemize}
\item \textbf{Decision Tree}: เป็นอัลกอริทึมที่สามารถอธิบายกฎการตัดสินใจได้ค่อนข้างชัดเจน จึงเหมาะสำหรับการพิจารณาโครงสร้างการตัดสินใจของการจำแนกผลลัพธ์

\item \textbf{Random Forest}: เป็นวิธีการรวมแบบจำลองต้นไม้ตัดสินใจหลายต้นเพื่อเพิ่มความแม่นยำ และช่วยลดปัญหาการเรียนรู้มากเกินไป (Overfitting)

\item \textbf{Support Vector Machine (SVM)}: ใช้สำหรับการจำแนกข้อมูลโดยอาศัยเส้นแบ่ง (Hyperplane) ซึ่งเหมาะกับข้อมูลที่มีความซับซ้อนและมีมิติสูง
\end{itemize}

\subsection{การฝึกสอนและทดสอบโมเดล (Training and Testing)}

ผู้วิจัยแบ่งชุดข้อมูลเป็นชุดฝึกสอนและชุดทดสอบ (Train/Test Split) ในสัดส่วนประมาณ 70--80/20--30 และใช้เทคนิค K-Fold Cross Validation เช่น 5-fold หรือ 10-fold เพื่อประเมินความเสถียรของโมเดล

\subsection{การประเมินโมเดล (Model Evaluation)}

ประเมินประสิทธิภาพของโมเดลด้วยตัวชี้วัด ได้แก่ Accuracy, Precision, Recall และ F1-score

\subsection{การอธิบายผลด้วย SHAP (Explainable AI)}

เพื่อเพิ่มความสามารถในการอธิบายผลของแบบจำลอง ผู้วิจัยใช้เทคนิค SHAP (SHapley Additive exPlanations) เพื่อวิเคราะห์ความสำคัญของตัวแปร และอธิบายอิทธิพลของตัวแปรที่มีต่อผลลัพธ์ทั้งในระดับภาพรวม (Global Explanation) และระดับรายบุคคล (Local Explanation) โดยผลลัพธ์ดังกล่าวนำไปใช้ประกอบการอภิปรายผลและจัดทำข้อเสนอเชิงนโยบายที่เหมาะสม

\subsection{การจำแนกกลุ่มเกษตรกร (Clustering Analysis)}

ผู้วิจัยใช้การวิเคราะห์แบบไม่กำกับ (Unsupervised Learning) เพื่อจัดกลุ่มเกษตรกรตามความคล้ายคลึงของคุณลักษณะ โดยประยุกต์ใช้วิธี K-Means Clustering เพื่อจำแนกกลุ่มจากชุดตัวแปรด้านศักยภาพ ทรัพยากร และทักษะ ทั้งนี้ ผลการจัดกลุ่มช่วยสนับสนุนการกำหนดแนวทางการพัฒนาทักษะและมาตรการสนับสนุนที่เหมาะสมแตกต่างกันตามลักษณะของแต่ละกลุ่มเป้าหมาย